% Generated by tools/build_v6.py; edit the assembly rules, not this file.
This section defines a scalar characterization that separates the part of
the law handled by rare-event scale entropy from a part handled by the
ordinary inverse-gap cost and a part left unresolved at the requested
accuracy. Its observable upper bound is proved in
Appendix~\ref{sec:entropy-tail-upper} and its matching lower bound in
Appendix~\ref{sec:matching-lower}. The decomposition is analytical:
the algorithm never observes a coin's component label.

\subsection{A homogeneous entropy and a selection functional}

For a finite submeasure $\sigma\le\mu$, use the canonical bands
$I_j=(b_j/2,b_j]$, $b_j=2^{-(j+1)}$, and put
\[
 a_j(\sigma)=\int_{I_j}d^{-2}\,\sigma(dd),\qquad
 A_\sigma=\sum_j a_j(\sigma).
\]
For a finite nonnegative vector $a$, define
\[
 \mathcal G(a)=A+\sum_j a_j\log(A/a_j),\qquad A=\sum_j a_j,
 \qquad \mathcal G(0)=0.
\]
Zero coordinates contribute zero. The functional $\mathcal G$ is concave, coordinatewise increasing,
homogeneous of degree one, and superadditive. Indeed, on positive
coordinates its derivative is $1+\log(A/a_j)$ and its Hessian is
$A^{-1}\mathbf1\mathbf1^\top-\operatorname{diag}(a_j^{-1})$,
which is negative semidefinite by Cauchy--Schwarz. Continuity gives the
boundary cases, and concavity plus homogeneity gives superadditivity. Define
\begin{equation}\label{eq:submeasure-entropy}
 \Gamma_\sigma=\sup_J\mathcal G(a_0(\sigma),\ldots,a_J(\sigma)).
\end{equation}
Thus $\Gamma_0=0$, and, when $0<A_\sigma<\infty$,
$\Gamma_\sigma=A_\sigma[1+\Ent((a_j/A_\sigma)_j)]$, with extended-value
interpretation. Also write
$H_\tau=\int d^{-1}\,\tau(dd)$ for any finite submeasure $\tau$.

Let $h=B^{-1/2}$ and define
\begin{equation}\label{eq:scalar-selection}
 \boxed{\quad
 \mathcal K_{\mu,\rho}(B)=
 \min\left\{\rho,\inf_{0\le\sigma\le\mu}
 \left[h\sqrt{\rho\Gamma_\sigma}
       +\int\min\{1,h/d\}\,d(\mu-\sigma)\right]\right\}.
 \quad}
\end{equation}
Equivalently,
\begin{equation}\label{eq:three-piece-functional}
 \mathcal K_{\mu,\rho}(B)=
 \min\left\{\rho,\inf_{\sigma+\tau+\zeta=\mu}
 \left[\zeta((0,1/2])
       +\frac{\sqrt{\rho\Gamma_\sigma}+H_\tau}{\sqrt B}\right]\right\}.
\end{equation}
Indeed, at fixed $\sigma$, assign each remaining magnitude to the cheaper
of the costs $1$ and $h/d$.

\begin{lemma}[Hard assignments suffice]\label{lem:hard-selection}
The infimum in \eqref{eq:scalar-selection} is unchanged if
$\sigma=\mu|S$ for a measurable set $S$. In particular, for a finite
atomic law a minimizing assignment can put each atom entirely in one of
the three pieces in \eqref{eq:three-piece-functional}.
\end{lemma}

\begin{proof}
Write $f=d\sigma/d\mu\in[0,1]$ and $S_t=\{f>t\}$, so
$\sigma=\int_0^1\mu|S_t\,dt$. Concavity of the square root of every
finite-fragment entropy gives
\[
 \sqrt{\Gamma_\sigma}\ge
       \int_0^1\sqrt{\Gamma_{\mu|S_t}}\,dt
\]
after monotone convergence in the fragment size. For a finite-objective
input the relevant entropies are finite for $t>0$, since
$\mu|S_t\le\sigma/t$ and
$\Gamma_{\mu|S_t}\le\Gamma_\sigma/t$. The remainder integral is linear,
so the objective at $\sigma$ is at least the average of the hard-set
objectives. Their infimum is therefore no larger. The reverse inequality
is immediate. For a finite law the minimum is attained among the finitely
many assignments. This is a selection problem, not a claim of convex
minimization over its assignments.
\end{proof}

The definition immediately implies
\[
 \mathcal K_{\mu,\rho}(B)\le
 \min\{\rho,\omega_\mu(B^{-1/2}),
                  \sqrt{\rho\Gamma_\mu/B}\},
 \qquad
 \omega_\mu(h)=\int\min\{1,h/d\}\,\mu(dd).
\]
Infinite quantities are allowed in this display.

\begin{theorem}[An entropy--tail law for the full joint problem]
\label{thm:entropy-tail-law}
There are universal constants $0<c<C<\infty$ such that, for every
probability measure $\mu$ on $(0,1/2]$, every $0<\rho\le1/4$, and every
$B>0$,
\begin{equation}\label{eq:entropy-tail-law}
 \boxed{\qquad
 c\,\mathcal K_{\mu,\rho}(B)\le R_{\mu,\rho}(B)
             \le C\,\mathcal K_{\mu,\rho}(B).
 \qquad}
\end{equation}
The functional has the equivalent selection and three-piece forms
\eqref{eq:scalar-selection}--\eqref{eq:three-piece-functional}, and hard
measurable assignments suffice. No density, regular variation, finite
inverse moment, or finite full-law entropy is required.
\end{theorem}

The upper half is Theorem~\ref{thm:three-piece-upper}; the lower half is
Theorem~\ref{thm:entropy-tail-lower}. The comparison constants in
\eqref{eq:entropy-tail-law} are universal but conservative. They are not
the factor sixteen proved for the finite convex program; these are
different representations and different comparison bounds.
The proof also gives the explicit conservative constants in
\eqref{eq:explicit-main-constants}.

\begin{theorem}[Three-piece observable upper bound]\label{thm:three-piece-upper}
There is a universal constant $C$ such that, for every magnitude law on
$(0,1/2]$, every $0<\rho\le1/4$, and every $B>0$,
\begin{equation}\label{eq:three-piece-upper}
 R_{\mu,\rho}(B)\le C\mathcal K_{\mu,\rho}(B).
\end{equation}
No component label or magnitude is revealed to the algorithm.
\end{theorem}
